\documentclass[11pt]{article}
\usepackage[margin=1in]{geometry}
\usepackage{amsmath,amssymb}
\usepackage{hyperref}

\newcommand{\E}{\mathbb E}
\newcommand{\cH}{\mathcal H}
\newcommand{\hbarp}{\hbar_{p,X}}
\newcommand{\betap}{\beta_{p,X}}

\title{Adapted antisymmetric stochastic integrals: the Hilbert-transform constant}
\author{arXiv automatic math run \texttt{fa\_banach\_001}}
\date{2026-06-18}

\begin{document}
\maketitle

\section*{Status}

This is a \emph{literature-implied answer (partial constant)} for the
antisymmetric route in Remark 5.12 of Yaroslavtsev's
\emph{Fourier multipliers and weak differential subordination of martingales in
UMD Banach spaces}, arXiv:1703.07817v5 \cite{Y17}.

The source remark asks for the analogue of Theorem 5.10 when the operator
\(A\) on the Hilbert noise space is antisymmetric.  Such an estimate with a
constant linear in \(\betap\) would imply the linear Hilbert-transform estimate
\[
\|\cH_X\|_{L^p(\mathbb R;X)\to L^p(\mathbb R;X)}
\le C_p\betap .
\]

\section*{Identification}

Let \(W^H\) be an \(H\)-cylindrical Brownian motion, let \(\Phi\) be an adapted
stochastic-integrable process with values in \(\mathcal L(H,X)\), and let
\[
 M=(\Phi\cdot W^H), \qquad N=((\Phi A)\cdot W^H),
\]
where \(A^*=-A\) and \(\|A\|\le 1\).  For every \(x^*\in X^*\),
\[
[\langle M,x^*\rangle,\langle N,x^*\rangle]_t
=\int_0^t\langle \Phi(s)^*x^*,A\Phi(s)^*x^*\rangle_H\,ds=0,
\]
so \(M\) and \(N\) are orthogonal in the sense of
Osekowski--Yaroslavtsev \cite{OY18}.  Also
\[
[\langle N,x^*\rangle]_t
=\int_0^t\|A\Phi(s)^*x^*\|_H^2\,ds
\le \int_0^t\|\Phi(s)^*x^*\|_H^2\,ds
=[\langle M,x^*\rangle]_t,
\]
so \(N\) is weakly differentially subordinate to \(M\).

Theorem 2.1 of \cite{OY18}, applied with
\(\Phi_0(x)=\Psi_0(x)=\|x\|^p\), gives
\[
\E\|N_t\|^p\le \hbarp^p \E\|M_t\|^p,\qquad t\ge 0,
\]
where \(\hbarp\) is the norm of the Hilbert transform on
\(L^p(\mathbb R;X)\).  The same consequence is restated in the introduction of
\cite{Y18SO}: orthogonal weakly differentially subordinate martingales satisfy
the sharp \(\hbarp\)-bound.

\section*{Scope and obstruction}

This does \emph{not} fully solve Remark 5.12 in the beta-constant form
suggested by Theorem 5.10.  It proves the adapted antisymmetric estimate with
\(\hbarp\), not with \(\betap\).  Osekowski--Yaroslavtsev explicitly record
that the case \(\Phi_0=\Psi_0=\|\cdot\|^p\) is tied to the celebrated open
problem whether \(\hbarp\lesssim_p \betap\); the best general upper bound
available there is still quadratic, \(\hbarp\le \betap^2\).

Therefore the current status is:
\[
\|((\Phi A)\cdot W^H)_t\|_{L^p(\Omega;X)}
\le \hbarp\,\|(\Phi\cdot W^H)_t\|_{L^p(\Omega;X)}
\]
for adapted antisymmetric \(A\), while the desired promotion to
\[
\|((\Phi A)\cdot W^H)_t\|_{L^p}
\le C_p\betap\,\|(\Phi\cdot W^H)_t\|_{L^p}
\]
would require resolving or bypassing the remaining Hilbert-transform/UMD
constant comparison problem in this route.

\begin{thebibliography}{9}

\bibitem{Y17}
I. Yaroslavtsev.
\newblock \emph{Fourier multipliers and weak differential subordination of
martingales in UMD Banach spaces}.
\newblock arXiv:1703.07817v5, 2017.

\bibitem{OY18}
A. Osekowski and I. Yaroslavtsev.
\newblock \emph{The Hilbert transform and orthogonal martingales in Banach
spaces}.
\newblock arXiv:1805.03948, 2018.

\bibitem{Y18SO}
I. Yaroslavtsev.
\newblock \emph{On strongly orthogonal martingales in UMD Banach spaces}.
\newblock arXiv:1812.08049, 2018.

\end{thebibliography}

\end{document}
